% Abstract.
%
% Every number here appears in a generated table under tables/. Two things this
% must not do, both decided and both mistakes the previous version made:
% promise a policy-impact metric (withdrawn, D18) or a single aggregate score
% (PF and RD are reported separately, D10). A third, learned from R18: do not
% sell d_trans -- the amendment is "report a measured distance", not "use ours".

\begin{abstract}
World models learn the dynamics of an environment so an agent can plan inside
them. When such a model is trained on one environment and then another, the
transition component that carries latent state forward is overwritten, and no
benchmark measures that degradation at the level of the component itself:
existing continual learning suites evaluate policies, or world models as
integrated systems. We build one --- three environment families, three levels
of dynamic distance each, five continual learning methods, five seeds, 225 runs
with 45 from-scratch reference pairs --- and report what it says about the way
this kind of experiment is usually designed and read.

Two results, both negative, and neither about which method wins. First, the
distance axis we built the benchmark around does not order forgetting.
Forgetting peaks at the \emph{medium} level in all three families, so the
labelled order is wrong in every one of them (Spearman $\rho = 0.05$ against
the labels; $0.58$ against how hard the new task is to fit). The cleanest case
needs no cross-family comparison: one family's maximum level is its medium
perturbation plus two more, applied to the same task pair, and it forgets
\emph{less}. Rerunning that pair at twice the training budget preserves the
ordering and refutes the obvious explanation --- the harder-perturbed task is
\emph{easier} to fit, not harder, because a cheetah at triple mass barely
moves. Forgetting tracked how much the new task demanded of the model, not how
far it sat from the old one.

Second, the forgetting happens almost entirely in the encoder, where
component-level metrics --- including ours --- cannot see it. Fine-tuning's
held-out reconstruction of the first task degrades by a factor of $745$ while
its prediction fidelity reads $-6.21$, i.e.\ improvement. EWC makes the
dissociation exact and predictable: its Fisher information is identically zero
on every encoder parameter, so it preserves the transition component almost
perfectly ($-0.07$) with an encoder degraded by $753$. A benchmark reporting
only the isolated metric would certify that these models did not forget.

We characterise five methods rather than ranking them, our own included, and
recommend that continual learning suites report a measured distance for each
task pair and hold it to account rather than treating level labels as an
experimental factor --- a standard we also apply to our own measured distance,
which fails it. Code, configurations and all 225 result files are released.
\end{abstract}
